\section*{Chapitre \ref{ch:b2:acyl-substitution} --- Dérivés des acides carboxyliques}

\begin{solution}{exo:b2:acyl-substitution:1}
Éthanamide $<$ éthanoate d'éthyle $<$ anhydride éthanoïque $<$ chlorure d'éthanoyle.
\end{solution}

\begin{solution}{exo:b2:acyl-substitution:2}
Acide éthanoïque (+ \ce{HCl})~; éthanoate d'éthyle~; éthanamide (+ chlorure d'ammonium, le second
équivalent d'ammoniac piégeant \ce{HCl})~; anhydride éthanoïque (+ \ce{NaCl})~;
\textit{N},\textit{N}-diméthyléthanamide (+ chlorure de triéthylammonium).
\end{solution}

\begin{solution}{exo:b2:acyl-substitution:3}
\ce{C6H5COOCH3 + NaOH -> C6H5COONa + CH3OH}. $13.6/136.15 = \qty{0.0999}{mol}$, autant de
\ce{NaOH}~: \qty{4.00}{g}.
\end{solution}

\begin{solution}{exo:b2:acyl-substitution:4}
$\gamma$-butyrolactone (oxolan-2-one)~: un cycle à cinq atomes, quatre carbones et un oxygène.
\end{solution}

\begin{solution}{exo:b2:acyl-substitution:5}
L'ammoniac s'additionne sur le carbone du carbonyle~; un proton passe de l'azote à l'oxygène (ou au
groupe partant)~; l'éthanolate (sous forme d'éthanol, après transfert de proton) part et la liaison
\ce{C=O} se reforme~: éthanamide et éthanol.
\end{solution}

\begin{solution}{exo:b2:acyl-substitution:6}
Le \ce{OH} phénolique est acylé (catalyse acide)~: acide 2-acétoxybenzoïque~; sous-produit~: acide éthanoïque.
\end{solution}

\begin{solution}{exo:b2:acyl-substitution:7}
Le 2-méthylpropan-2-ol. La première addition donne la propanone, plus électrophile que l'ester, qui
capte aussitôt le second méthyle.
\end{solution}

\begin{solution}{exo:b2:acyl-substitution:8}
Le butane-1,4-diol~: l'ester est réduit en deux alcools primaires, le cycle étant ouvert.
\end{solution}

\begin{solution}{exo:b2:acyl-substitution:9}
Un équivalent~: $x^2/(1 - x)^2 = 4$, $x = 2/3$. Trois équivalents~: $x^2/[(1 - x)(3 - x)] = 4$, $3x^2 - 16x
+ 12 = 0$, $x = 0.90$. Éliminer l'eau au fur et à mesure de sa formation (Dean--Stark) déplace encore la réaction.
\end{solution}

\begin{solution}{exo:b2:acyl-substitution:10}
\ce{SOCl2} donne le chlorure de benzoyle~; avec deux équivalents de diéthylamine (ou un seul plus
de la triéthylamine), il donne l'amide. L'acide et l'amine mélangés directement subissent une réaction
acide--base~: on obtient le benzoate de diéthylammonium, qui ne donne l'amide que par chauffage intense.
\end{solution}

\begin{solution}{exo:b2:acyl-substitution:11}
Trois KOH par trioléine~: $3 \times 56.11/885.4 = \qty{0.190}{g}$ par gramme, soit un indice de
\omterm{def:b2:acyl-substitution:saponification}{saponification} de 190.
\end{solution}

\begin{solution}{exo:b2:acyl-substitution:12}
\ce{NaBH4}~: seulement la cétone, en alcool secondaire. \ce{LiAlH4}~: la cétone en alcool secondaire,
l'ester en alcool primaire (en libérant la partie alcool), l'amide en amine.
\end{solution}

\begin{solution}{pb:b2:acyl-substitution:1}
\textbf{1.} Acide oléique \ce{C18H34O2}~; glycérol \ce{C3H8O3}.
\textbf{2.} \ce{C57H104O6}~: $57 \times 12.011 + 104 \times 1.008 + 6 \times 15.999 = \qty{885.45}{g/mol}$.
\textbf{3.} Trois.
\textbf{4.} La double liaison cis coude les chaînes, qui s'empilent mal~: forces intermoléculaires
faibles, température de fusion basse.
\textbf{5.} Du glycérol et trois oléates de sodium, un savon.
\textbf{6.} Un mélange d'esters méthyliques d'acides gras.
\textbf{7.} \ce{C57H104O6 + 3CH3OH -> C3H8O3 + 3C19H36O2}.
\textbf{8.} L'ion méthanolate, formé à partir du méthanol et de la base, est un nucléophile bien plus
fort que le méthanol~; il attaque les carbonyles des esters.
\textbf{9.} Le carbone du carbonyle d'un groupe ester, portant \ce{O-}, le \ce{OCH3} du méthanolate et
l'oxygène du glycéryle~: tétraédrique.
\textbf{10.} La \omterm{def:b2:acyl-substitution:transesterification}{transestérification} est un équilibre de constante voisine de 1~; le méthanol en excès le
déplace vers les esters méthyliques.
\textbf{11.} Le glycérol se sépare en une phase inférieure~: retirer un produit déplace l'équilibre.
\textbf{12.} Non~: le méthanolate est régénéré à chaque échange.
\textbf{13.} Il neutralise la base, en donnant de l'oléate de sodium (un savon)~: le catalyseur est consommé.
\textbf{14.} Elle les saponifie~: l'hydroxyde (issu de l'eau et du méthanolate) coupe les esters en savon
et méthanol, de façon irréversible.
\textbf{15.} Il consomme de la base, émulsionne le mélange et empêche la phase de glycérol de se séparer.
\textbf{16.} Par une \omterm{def:b2:acyl-substitution:esterification}{estérification} préalable par le méthanol en catalyse acide, qui transforme les acides
libres en esters méthyliques.
\textbf{17.} L'eau hydrolyse les esters et transforme le catalyseur en hydroxyde, qui saponifie.
\textbf{18.} $10^6/885.45 = \qty{1129}{mol}$.
\textbf{19.} $3 \times 1129 \times 32.04 = \qty{108.6}{kg}$.
\textbf{20.} $3 \times 1129 \times 296.50 = \qty{1004.6}{kg}$.
\textbf{21.} $1000 + 108.6 = \qty{1108.6}{kg}$ en entrée~; $104.0 + 1004.6 = \qty{1108.6}{kg}$ en sortie.
\textbf{22.} \qty{20}{kg} d'acide oléique, $20\,000/282.47 = \qty{70.8}{mol}$, autant de \ce{NaOH}~:
\qty{2.83}{kg}.
\textbf{23.} $1129 \times 92.09 = \boldsymbol{\approx \qty{104}{kg}}$ de glycérol par tonne.
\end{solution}
